Table~\ref{tab:error-structure} showed distinct error regimes across the four conditions based on the mean per-session false positives (FP), false negatives (FN), and FP:FN ratio. BLINKER-concat was strongly false-positive-heavy, with 641.5 FP, 20.4 FN, and an FP:FN ratio of 31.4, reflecting severe over-detection. MNE-annot was false-negative-heavy, with 207.6 FP, 531.8 FN, and a ratio of 0.39, and Proposed-Mean was mildly false-negative-heavy, with 162.6 FP, 184.4 FN, and a ratio of 0.88. In contrast, Proposed-Med produced the most balanced error profile, with 180.7 FP, 151.7 FN, and a near-symmetric ratio of 1.19. Overall, Proposed-Med occupied a balanced operating point between the strongly over-detecting BLINKER-concat and the under-detecting MNE-annot, dramatically reducing the false-positive burden relative to BLINKER-concat while keeping false negatives close to false positives.
